% Part: normal-modal-logic
% Chapter: frame-correspondence
% Section: equivalence-S5

\documentclass[../../../include/open-logic-section]{subfiles}

\begin{document}

\olfileid[pt]{nml}{frd}{es5}

\olsection{Relações de equivalência e \Log{S5}}

A lógica modal \Log{S5} é caracterizada como o conjunto de !!{formula}s
válidas em todos os frames universais, O.s.\@, todo mundo é acessível a partir
de todo mundo, inclusive de si mesmo. Nesse cenário, $\Box$ corresponde à
necessidade e $\Diamond$ à possibilidade: $\Box !A$ é verdadeira se $!A$ é
verdadeira em \emph{todo} mundo, e $\Diamond !A$ é verdadeira se $!A$ é
verdadeira em \emph{algum} mundo. Ocorre que \Log{S5} também pode ser
caracterizada como as !!{formula}s válidas em todos os frames reflexivos,
simétricos e transitivos, O.s.\@, em todas as \emph{relações de equivalência}.

\begin{defn}
  Uma relação binária $R$ sobre $W$ é uma \emph{relação de equivalência} se e
  somente se é reflexiva, simétrica e transitiva. Uma relação $R$ sobre $W$ é
  \emph{universal} se e somente se $Ruv$ para todos $u,v \in W$.
\end{defn}

Como \Ax{T}, \Ax{B} e \Ax{4} caracterizam, respectivamente, os frames
reflexivos, simétricos e transitivos, os frames nos quais a relação de
acessibilidade é uma relação de equivalência são exatamente aqueles nos quais
as três !!{formula}s são válidas. Ocorre que as relações de equivalência
também podem ser caracterizadas por outras combinações de !!{formula}s, pois
as condições com as quais definimos as relações de equivalência são
equivalentes a combinações de outras condições conhecidas sobre~$R$.

\begin{prop}\ollabel{prop:equivalences}
  As seguintes condições são equivalentes:
  \begin{enumerate}
  \item $R$ é uma relação de equivalência;
  \item $R$ é reflexiva e euclidiana;
  \item $R$ é serial, simétrica e euclidiana;
  \item $R$ é serial, simétrica e transitiva.
  \end{enumerate}
\end{prop}

\begin{proof}
  Exercício.
\end{proof}

\begin{prob}
  Demonstre \olref[nml][frd][es5]{prop:equivalences} mostrando:
  \begin{enumerate}
  \item Se $R$ é simétrica e transitiva, então é euclidiana.
  \item Se $R$ é reflexiva, então é serial.
  \item Se $R$ é reflexiva e euclidiana, então é simétrica.
  \item Se $R$ é simétrica e euclidiana, então é transitiva.
  \item Se $R$ é serial, simétrica e transitiva, então é reflexiva.
  \end{enumerate}
  Explique por que isso basta para demonstrar que as condições são
  equivalentes.
\end{prob}

\olref{prop:equivalences} é a contraparte semântica de
\olref[prf][prs]{prop:S5}, pois fornece uma caracterização equivalente da
lógica modal dos frames sobre os quais $R$ é uma relação de equivalência (a
lógica tradicionalmente chamada de \Log{S5}).

Qual é a relação entre relações universais e relações de equivalência? Embora
toda relação universal seja uma relação de equivalência, é claro que nem toda
relação de equivalência é universal. Contudo, as !!{formula}s válidas em todas
as relações universais são exatamente as mesmas que são válidas em todas as
relações de equivalência.

\begin{prop}
  Seja $R$ uma relação de equivalência e, para cada $w \in W$, defina a
  \emph{classe de equivalência} de $w$ como o conjunto
  $[w] = \{w'\in W : Rww'\}$. Então:
  \begin{enumerate}
  \item $w \in [w]$;
  \item $R$ é universal em cada classe de equivalência $[w]$;
  \item A coleção de classes de equivalência particiona $W$ em subconjuntos
    mutuamente exclusivos e conjuntamente exaustivos.
  \end{enumerate}
\end{prop}

\begin{prop}\ollabel{prop:S5=univ}
  !!^a{formula}~$!A$ é válida em todos os frames
  $\mModel{F} =\tuple{W,R}$ nos quais $R$ é uma relação de equivalência se e
  somente se é válida em todos os frames $\mModel{F} =\tuple{W,R}$ nos quais
  $R$ é universal. Portanto, a lógica dos frames universais é simplesmente
  \Log{S5}.
\end{prop}

\begin{proof}
  É imediato verificar que uma relação universal~$R$ sobre~$W$ é uma relação
  de equivalência. Portanto, se $!A$ é válida em todos os frames nos quais
  $R$ é uma relação de equivalência, então é válida em todos os frames
  universais. Para a outra direção, argumentamos contrapositivamente: suponha
  que $!B$ seja !!a{formula} que falha em um mundo $w$ de um modelo
  $\mModel{M} = \tuple{W, R, V}$ baseado em um frame $\tuple{W,R}$, no qual
  $R$ é uma relação de equivalência sobre $W$. Assim,
  $\mSat/{M}{!B}[w]$. Defina um modelo
  $\mModel{M}' = \tuple{W', R', V'}$ da seguinte maneira:
  \begin{enumerate}
  \item $W' = [w]$;
  \item $R'$ é universal sobre $W'$;
  \item $V'(p) = V(p) \cap W'$.
  \end{enumerate}
  (Assim, o conjunto $W'$ de mundos em $\mModel{M}'$ é representado pela
  área sombreada em \olref{fig:partition}.) É fácil ver que $R$ e $R'$
  coincidem em $W'$. Pode-se então mostrar, por indução sobre as !!{formula}s,
  que, para todo $w' \in W'$: $\mSat{M'}{!A}[w']$ se e somente se
  $\mSat{M}{!A}[w']$ para cada $!A$ (isso faz sentido, pois $W' \subseteq W$).
  Em particular, $\mSat/{M'}{!B}[w]$, e $!B$ falha em um modelo baseado em um
  frame universal.
\end{proof}

\begin{figure}[t]
  \centering
  \begin{tikzpicture}[node distance=2cm, auto, thick]
    \clip [rounded corners] (0,0) -- (6,0) -- (6,4) -- (0,4) --  cycle;
    \begin{scope}
      \clip (0,0) -- (2,0)  .. controls (1.5,1.5) and (4.5,1.5)
      .. (4,4) -- (0,4) -- (0,0) -- cycle;
      \filldraw[fill=gray!40] (0,4) -- (0,2) .. controls (1.5,1.5)
      and (4.5,1.5) .. (4.5,0) -- (6,0) -- (6,4) -- (0,4) --  cycle;
    \end{scope}
    \draw [name path=line1] (2,0) .. controls (1.5,1.5) and (4.5,1.5) .. (4,4);
    \draw [name path=line2] (0,2)  .. controls (1.5,1.5) and (4.5,1.5) .. (4.5,0);
    \path [name intersections={of=line1 and line2}];
    \draw [rounded corners, use as bounding box, very thick] (0,0)
    -- (6,0) -- (6,4) -- (0,4) --  cycle; 
    \node at (2,3) {$[w]$} ;
    \node at (1,1) {$[u]$} ;
    \node at (3,0.75) {$[v]$} ;
    \node at (4.75,2) {$[z]$} ;
  \end{tikzpicture}
  \caption{Uma partição de $W$ em classes de equivalência.}
  \ollabel{fig:partition}
\end{figure}

\end{document}

